{\scriptsize\tiny
    \begin{minipage}[t]{0.54\textwidth}
        \begin{tabularx}{\linewidth}{|p{0.35\textwidth}|X|}
            \hline 
            \textbf{Notions et contenus} & \textbf{Capacités exigibles} \\
            \hline 
            \multicolumn{2}{|l|}{5.2. Électrostatique } \\
            \hline 
            \multicolumn{2}{|l|}{5.2.1. Champ électrostatique} \\
            \hline 
            Loi de Coulomb.
            Champ et potentiel électrostatiques créés par une charge ponctuelle. Principe de superposition. & 
            Exprimer le champ électrostatique et le potentiel créés par une distribution discrète de charges.
            Citer quelques ordres de grandeur de champs électrostatiques. \\
            \hline 
            Propriétés du champ électrostatique
Symétries. & Exploiter les propriétés de symétrie des sources (translation, rotation, symétrie plane, conjugaison de charges) pour prévoir des propriétés du champ créé. \\
            \hline 
            Circulation du champ électrostatique. Potentiel électrostatique. Équations locales. & Relier l'existence d'un potentiel électrostatique à la nullité du rotationnel du vecteur champ électrostatique. Justifier l'orthogonalité des lignes de champ avec les surfaces équipotentielles et leur orientation dans le sens des potentiels décroissants.\\
            \hline
        \end{tabularx}
\end{minipage}
\begin{minipage}[t]{0.43\textwidth}
    \begin{tabularx}{\linewidth}{|p{0.35\textwidth}|X|}
        \hline 
        Théorème de Gauss et équation locale de Maxwell-Gauss. & Choisir une surface adaptée et utiliser le théorème de Gauss. \\
        \hline 
        Lignes de champ électrostatique. Équipotentielles. & Justifier qu'une carte de lignes de champ puisse ou non être celle d'un champ électrostatique.
        Repérer, sur une carte de champ électrostatique, d'éventuelles sources du champ et leur signe.
        Associer l'évolution de la norme du champ électrostatique à l'évasement des tubes de champ loin des sources.
        Relier équipotentielles et lignes de champ électrostatique.
        Évaluer la norme du champ électrostatique à partir d'un réseau de lignes équipotentielles. \\
        \hline
        \multicolumn{2}{|l|}{ 5.2.3. Analogies avec le champ gravitationnel } \\ 
        \hline
        Analogies entre champ électrostatique et champ gravitationnel. & Utiliser les analogies entre les forces électrostatique et gravitationnelle pour determiner l'expression de champs gravitationnels. \\ 
        \hline
    \end{tabularx}
\end{minipage}
}
